\section{余下账本事件导航：western-han}

\label{sec:remaining-event-anchors-western-han}

本节为尚未在正文设置目标的账本记录建立直接跳转。锚点显示可核日期与事件名称；其解释仍应回到账本所列来源、置信提示及既有正文的区域和材料边界。

\subsection{事件导航与来源边界}

\eventanchor{WH-LEDGER-004-2}{前207年冬·“章邯向项羽降，秦军降卒后遭坑杀的叙事引发争议”的命令、联盟或地方响应}

\eventanchor{WH-LEDGER-004-3}{前207年冬·章邯向项羽降，秦军降卒后遭坑杀的叙事引发争议}

\eventanchor{WH-LEDGER-004-4}{前207年冬·“章邯向项羽降，秦军降卒后遭坑杀的叙事引发争议”的执行中的空间与人员处置}

\eventanchor{WH-LEDGER-004-5}{前207年冬·“章邯向项羽降，秦军降卒后遭坑杀的叙事引发争议”的后续制度或军政调整}

\eventanchor{WH-LEDGER-004-6}{前207年冬·“章邯向项羽降，秦军降卒后遭坑杀的叙事引发争议”的异源材料与影响范围的复核}

\eventanchor{WH-LEDGER-005-1}{前207年十月·“刘邦军先入关中，子婴降于轵道”的前期动员与资源准备}

\eventanchor{WH-LEDGER-005-2}{前207年十月·“刘邦军先入关中，子婴降于轵道”的命令、联盟或地方响应}

\eventanchor{WH-LEDGER-005-3}{前207年十月·刘邦军先入关中，子婴降于轵道}

\eventanchor{WH-LEDGER-005-4}{前207年十月·“刘邦军先入关中，子婴降于轵道”的执行中的空间与人员处置}

\eventanchor{WH-LEDGER-005-5}{前207年十月·“刘邦军先入关中，子婴降于轵道”的后续制度或军政调整}

\eventanchor{WH-LEDGER-005-6}{前207年十月·“刘邦军先入关中，子婴降于轵道”的异源材料与影响范围的复核}

\eventanchor{WH-LEDGER-006-1}{前206年·“鸿门会后项羽控制关中，刘邦暂避其锋”的前期动员与资源准备}

\eventanchor{WH-LEDGER-006-2}{前206年·“鸿门会后项羽控制关中，刘邦暂避其锋”的命令、联盟或地方响应}

\eventanchor{WH-LEDGER-006-3}{前206年·鸿门会后项羽控制关中，刘邦暂避其锋}

\eventanchor{WH-LEDGER-006-4}{前206年·“鸿门会后项羽控制关中，刘邦暂避其锋”的执行中的空间与人员处置}

\eventanchor{WH-LEDGER-006-5}{前206年·“鸿门会后项羽控制关中，刘邦暂避其锋”的后续制度或军政调整}

\eventanchor{WH-LEDGER-006-6}{前206年·“鸿门会后项羽控制关中，刘邦暂避其锋”的异源材料与影响范围的复核}

\eventanchor{WH-LEDGER-007-1}{前206年·“项羽分封十八王，自立西楚霸王，刘邦徙封汉中”的前期动员与资源准备}

\eventanchor{WH-LEDGER-007-2}{前206年·“项羽分封十八王，自立西楚霸王，刘邦徙封汉中”的命令、联盟或地方响应}

\eventanchor{WH-LEDGER-007-3}{前206年·项羽分封十八王，自立西楚霸王，刘邦徙封汉中}

\eventanchor{WH-LEDGER-007-4}{前206年·“项羽分封十八王，自立西楚霸王，刘邦徙封汉中”的执行中的空间与人员处置}

\eventanchor{WH-LEDGER-007-5}{前206年·“项羽分封十八王，自立西楚霸王，刘邦徙封汉中”的后续制度或军政调整}

\eventanchor{WH-LEDGER-007-6}{前206年·“项羽分封十八王，自立西楚霸王，刘邦徙封汉中”的异源材料与影响范围的复核}

\eventanchor{WH-LEDGER-008-1}{前206年八月至九月·“刘邦出汉中，平定雍、塞、翟三秦”的前期动员与资源准备}

\eventanchor{WH-LEDGER-008-2}{前206年八月至九月·“刘邦出汉中，平定雍、塞、翟三秦”的命令、联盟或地方响应}

\eventanchor{WH-LEDGER-008-3}{前206年八月至九月·刘邦出汉中，平定雍、塞、翟三秦}

\eventanchor{WH-LEDGER-008-4}{前206年八月至九月·“刘邦出汉中，平定雍、塞、翟三秦”的执行中的空间与人员处置}

\eventanchor{WH-LEDGER-008-5}{前206年八月至九月·“刘邦出汉中，平定雍、塞、翟三秦”的后续制度或军政调整}

\eventanchor{WH-LEDGER-008-6}{前206年八月至九月·“刘邦出汉中，平定雍、塞、翟三秦”的异源材料与影响范围的复核}

\eventanchor{WH-LEDGER-009-1}{前205年·“刘邦联军攻入彭城后遭项羽反击而溃败”的前期动员与资源准备}

\eventanchor{WH-LEDGER-009-2}{前205年·“刘邦联军攻入彭城后遭项羽反击而溃败”的命令、联盟或地方响应}

\eventanchor{WH-LEDGER-009-3}{前205年·刘邦联军攻入彭城后遭项羽反击而溃败}

\eventanchor{WH-LEDGER-009-4}{前205年·“刘邦联军攻入彭城后遭项羽反击而溃败”的执行中的空间与人员处置}

\eventanchor{WH-LEDGER-009-5}{前205年·“刘邦联军攻入彭城后遭项羽反击而溃败”的后续制度或军政调整}

\eventanchor{WH-LEDGER-009-6}{前205年·“刘邦联军攻入彭城后遭项羽反击而溃败”的异源材料与影响范围的复核}

\eventanchor{WH-LEDGER-010-1}{前205年·“韩信、曹参等取魏地，汉军夺取河东交通节点”的前期动员与资源准备}

\eventanchor{WH-LEDGER-010-2}{前205年·“韩信、曹参等取魏地，汉军夺取河东交通节点”的命令、联盟或地方响应}

\eventanchor{WH-LEDGER-010-3}{前205年·韩信、曹参等取魏地，汉军夺取河东交通节点}

\eventanchor{WH-LEDGER-010-4}{前205年·“韩信、曹参等取魏地，汉军夺取河东交通节点”的执行中的空间与人员处置}

\eventanchor{WH-LEDGER-010-5}{前205年·“韩信、曹参等取魏地，汉军夺取河东交通节点”的后续制度或军政调整}

\eventanchor{WH-LEDGER-010-6}{前205年·“韩信、曹参等取魏地，汉军夺取河东交通节点”的异源材料与影响范围的复核}

\eventanchor{WH-LEDGER-011-1}{前204年·“韩信破赵于井陉，收编赵地兵员”的前期动员与资源准备}

\eventanchor{WH-LEDGER-011-2}{前204年·“韩信破赵于井陉，收编赵地兵员”的命令、联盟或地方响应}

\eventanchor{WH-LEDGER-011-3}{前204年·韩信破赵于井陉，收编赵地兵员}

\eventanchor{WH-LEDGER-011-4}{前204年·“韩信破赵于井陉，收编赵地兵员”的执行中的空间与人员处置}

\eventanchor{WH-LEDGER-011-5}{前204年·“韩信破赵于井陉，收编赵地兵员”的后续制度或军政调整}

\eventanchor{WH-LEDGER-011-6}{前204年·“韩信破赵于井陉，收编赵地兵员”的异源材料与影响范围的复核}

\eventanchor{WH-LEDGER-012-1}{前204年·“荥阳、成皋争夺持续，刘邦一度突围西走”的前期动员与资源准备}

\eventanchor{WH-LEDGER-012-2}{前204年·“荥阳、成皋争夺持续，刘邦一度突围西走”的命令、联盟或地方响应}

\eventanchor{WH-LEDGER-012-3}{前204年·荥阳、成皋争夺持续，刘邦一度突围西走}

\eventanchor{WH-LEDGER-012-4}{前204年·“荥阳、成皋争夺持续，刘邦一度突围西走”的执行中的空间与人员处置}

\eventanchor{WH-LEDGER-012-5}{前204年·“荥阳、成皋争夺持续，刘邦一度突围西走”的后续制度或军政调整}

\eventanchor{WH-LEDGER-012-6}{前204年·“荥阳、成皋争夺持续，刘邦一度突围西走”的异源材料与影响范围的复核}

\eventanchor{WH-LEDGER-013-1}{前203年·“韩信击齐并在潍水一带破楚援军”的前期动员与资源准备}

\eventanchor{WH-LEDGER-013-2}{前203年·“韩信击齐并在潍水一带破楚援军”的命令、联盟或地方响应}

\eventanchor{WH-LEDGER-013-3}{前203年·韩信击齐并在潍水一带破楚援军}

\eventanchor{WH-LEDGER-013-4}{前203年·“韩信击齐并在潍水一带破楚援军”的执行中的空间与人员处置}

\eventanchor{WH-LEDGER-013-5}{前203年·“韩信击齐并在潍水一带破楚援军”的后续制度或军政调整}

\eventanchor{WH-LEDGER-013-6}{前203年·“韩信击齐并在潍水一带破楚援军”的异源材料与影响范围的复核}

\eventanchor{WH-LEDGER-014-1}{前203年·“楚汉以鸿沟为界议和，旋即再战”的前期动员与资源准备}

\subsection{事件导航与来源边界}

\eventanchor{WH-LEDGER-014-2}{前203年·“楚汉以鸿沟为界议和，旋即再战”的命令、联盟或地方响应}

\eventanchor{WH-LEDGER-014-3}{前203年·楚汉以鸿沟为界议和，旋即再战}

\eventanchor{WH-LEDGER-014-4}{前203年·“楚汉以鸿沟为界议和，旋即再战”的执行中的空间与人员处置}

\eventanchor{WH-LEDGER-014-5}{前203年·“楚汉以鸿沟为界议和，旋即再战”的后续制度或军政调整}

\eventanchor{WH-LEDGER-014-6}{前203年·“楚汉以鸿沟为界议和，旋即再战”的异源材料与影响范围的复核}

\eventanchor{WH-LEDGER-015-1}{前202年·“汉军及诸侯合围项羽于垓下，项羽败死”的前期动员与资源准备}

\eventanchor{WH-LEDGER-015-2}{前202年·“汉军及诸侯合围项羽于垓下，项羽败死”的命令、联盟或地方响应}

\eventanchor{WH-LEDGER-015-3}{前202年·汉军及诸侯合围项羽于垓下，项羽败死}

\eventanchor{WH-LEDGER-015-4}{前202年·“汉军及诸侯合围项羽于垓下，项羽败死”的执行中的空间与人员处置}

\eventanchor{WH-LEDGER-015-5}{前202年·“汉军及诸侯合围项羽于垓下，项羽败死”的后续制度或军政调整}

\eventanchor{WH-LEDGER-015-6}{前202年·“汉军及诸侯合围项羽于垓下，项羽败死”的异源材料与影响范围的复核}

\eventanchor{WH-LEDGER-016-1}{前202年二月·“刘邦即皇帝位，汉朝建立”的前期动员与资源准备}

\eventanchor{WH-LEDGER-016-2}{前202年二月·“刘邦即皇帝位，汉朝建立”的命令、联盟或地方响应}

\eventanchor{WH-LEDGER-016-3}{前202年二月·刘邦即皇帝位，汉朝建立}

\eventanchor{WH-LEDGER-016-4}{前202年二月·“刘邦即皇帝位，汉朝建立”的执行中的空间与人员处置}

\eventanchor{WH-LEDGER-016-5}{前202年二月·“刘邦即皇帝位，汉朝建立”的后续制度或军政调整}

\eventanchor{WH-LEDGER-016-6}{前202年二月·“刘邦即皇帝位，汉朝建立”的异源材料与影响范围的复核}

\eventanchor{WH-LEDGER-017-1}{前202年五月·“汉廷定都长安，依托关中经营全国”的前期动员与资源准备}

\eventanchor{WH-LEDGER-017-2}{前202年五月·“汉廷定都长安，依托关中经营全国”的命令、联盟或地方响应}

\eventanchor{WH-LEDGER-017-3}{前202年五月·汉廷定都长安，依托关中经营全国}

\eventanchor{WH-LEDGER-017-4}{前202年五月·“汉廷定都长安，依托关中经营全国”的执行中的空间与人员处置}

\eventanchor{WH-LEDGER-017-5}{前202年五月·“汉廷定都长安，依托关中经营全国”的后续制度或军政调整}

\eventanchor{WH-LEDGER-017-6}{前202年五月·“汉廷定都长安，依托关中经营全国”的异源材料与影响范围的复核}

\eventanchor{WH-LEDGER-018-1}{前201年·“汉廷大规模封建功臣与异姓诸侯王”的前期动员与资源准备}

\eventanchor{WH-LEDGER-018-2}{前201年·“汉廷大规模封建功臣与异姓诸侯王”的命令、联盟或地方响应}

\eventanchor{WH-LEDGER-018-3}{前201年·汉廷大规模封建功臣与异姓诸侯王}

\eventanchor{WH-LEDGER-018-4}{前201年·“汉廷大规模封建功臣与异姓诸侯王”的执行中的空间与人员处置}

\eventanchor{WH-LEDGER-018-5}{前201年·“汉廷大规模封建功臣与异姓诸侯王”的后续制度或军政调整}

\eventanchor{WH-LEDGER-018-6}{前201年·“汉廷大规模封建功臣与异姓诸侯王”的异源材料与影响范围的复核}

\eventanchor{WH-LEDGER-019-1}{前200年冬·“刘邦北征受困白登，与匈奴议和撤军”的前期动员与资源准备}

\eventanchor{WH-LEDGER-019-2}{前200年冬·“刘邦北征受困白登，与匈奴议和撤军”的命令、联盟或地方响应}

\eventanchor{WH-LEDGER-019-3}{前200年冬·刘邦北征受困白登，与匈奴议和撤军}

\eventanchor{WH-LEDGER-019-4}{前200年冬·“刘邦北征受困白登，与匈奴议和撤军”的执行中的空间与人员处置}

\eventanchor{WH-LEDGER-019-5}{前200年冬·“刘邦北征受困白登，与匈奴议和撤军”的后续制度或军政调整}

\eventanchor{WH-LEDGER-019-6}{前200年冬·“刘邦北征受困白登，与匈奴议和撤军”的异源材料与影响范围的复核}

\eventanchor{WH-LEDGER-020-1}{前198年·“韩信由楚王徙为淮阴侯，异姓王权受抑”的前期动员与资源准备}

\eventanchor{WH-LEDGER-020-2}{前198年·“韩信由楚王徙为淮阴侯，异姓王权受抑”的命令、联盟或地方响应}

\eventanchor{WH-LEDGER-020-3}{前198年·韩信由楚王徙为淮阴侯，异姓王权受抑}

\eventanchor{WH-LEDGER-020-4}{前198年·“韩信由楚王徙为淮阴侯，异姓王权受抑”的执行中的空间与人员处置}

\eventanchor{WH-LEDGER-020-5}{前198年·“韩信由楚王徙为淮阴侯，异姓王权受抑”的后续制度或军政调整}

\eventanchor{WH-LEDGER-020-6}{前198年·“韩信由楚王徙为淮阴侯，异姓王权受抑”的异源材料与影响范围的复核}

\eventanchor{WH-LEDGER-021-1}{前196年·“吕后与萧何诱杀韩信于长乐宫”的前期动员与资源准备}

\eventanchor{WH-LEDGER-021-2}{前196年·“吕后与萧何诱杀韩信于长乐宫”的命令、联盟或地方响应}

\eventanchor{WH-LEDGER-021-3}{前196年·吕后与萧何诱杀韩信于长乐宫}

\eventanchor{WH-LEDGER-021-4}{前196年·“吕后与萧何诱杀韩信于长乐宫”的执行中的空间与人员处置}

\eventanchor{WH-LEDGER-021-5}{前196年·“吕后与萧何诱杀韩信于长乐宫”的后续制度或军政调整}

\eventanchor{WH-LEDGER-021-6}{前196年·“吕后与萧何诱杀韩信于长乐宫”的异源材料与影响范围的复核}

\eventanchor{WH-LEDGER-022-1}{前196年前后·“刘敬建议迁徙关东豪族十余万口入关中以实京师”的前期动员与资源准备}

\eventanchor{WH-LEDGER-022-2}{前196年前后·“刘敬建议迁徙关东豪族十余万口入关中以实京师”的命令、联盟或地方响应}

\eventanchor{WH-LEDGER-022-3}{前196年前后·刘敬建议迁徙关东豪族十余万口入关中以实京师}

\eventanchor{WH-LEDGER-022-4}{前196年前后·“刘敬建议迁徙关东豪族十余万口入关中以实京师”的执行中的空间与人员处置}

\eventanchor{WH-LEDGER-022-5}{前196年前后·“刘敬建议迁徙关东豪族十余万口入关中以实京师”的后续制度或军政调整}

\eventanchor{WH-LEDGER-022-6}{前196年前后·“刘敬建议迁徙关东豪族十余万口入关中以实京师”的异源材料与影响范围的复核}

\eventanchor{WH-LEDGER-023-1}{前196年·“梁王彭越被废杀，梁地改由宗室控制”的前期动员与资源准备}

\eventanchor{WH-LEDGER-023-2}{前196年·“梁王彭越被废杀，梁地改由宗室控制”的命令、联盟或地方响应}

\eventanchor{WH-LEDGER-023-3}{前196年·梁王彭越被废杀，梁地改由宗室控制}

\eventanchor{WH-LEDGER-023-4}{前196年·“梁王彭越被废杀，梁地改由宗室控制”的执行中的空间与人员处置}

\eventanchor{WH-LEDGER-023-5}{前196年·“梁王彭越被废杀，梁地改由宗室控制”的后续制度或军政调整}

\eventanchor{WH-LEDGER-023-6}{前196年·“梁王彭越被废杀，梁地改由宗室控制”的异源材料与影响范围的复核}

\eventanchor{WH-LEDGER-024-1}{前195年·“淮南王英布反汉，刘邦亲征后病逝”的前期动员与资源准备}

\subsection{事件导航与来源边界}

\eventanchor{WH-LEDGER-024-2}{前195年·“淮南王英布反汉，刘邦亲征后病逝”的命令、联盟或地方响应}

\eventanchor{WH-LEDGER-024-3}{前195年·淮南王英布反汉，刘邦亲征后病逝}

\eventanchor{WH-LEDGER-024-4}{前195年·“淮南王英布反汉，刘邦亲征后病逝”的执行中的空间与人员处置}

\eventanchor{WH-LEDGER-024-5}{前195年·“淮南王英布反汉，刘邦亲征后病逝”的后续制度或军政调整}

\eventanchor{WH-LEDGER-024-6}{前195年·“淮南王英布反汉，刘邦亲征后病逝”的异源材料与影响范围的复核}

\eventanchor{WH-LEDGER-025-1}{前195年前后·“朝廷以白马盟誓限制非刘氏封王”的前期动员与资源准备}

\eventanchor{WH-LEDGER-025-2}{前195年前后·“朝廷以白马盟誓限制非刘氏封王”的命令、联盟或地方响应}

\eventanchor{WH-LEDGER-025-3}{前195年前后·朝廷以白马盟誓限制非刘氏封王}

\eventanchor{WH-LEDGER-025-4}{前195年前后·“朝廷以白马盟誓限制非刘氏封王”的执行中的空间与人员处置}

\eventanchor{WH-LEDGER-025-5}{前195年前后·“朝廷以白马盟誓限制非刘氏封王”的后续制度或军政调整}

\eventanchor{WH-LEDGER-025-6}{前195年前后·“朝廷以白马盟誓限制非刘氏封王”的异源材料与影响范围的复核}

\eventanchor{WH-LEDGER-026-1}{前195年四月·“高祖刘邦去世，惠帝即位，吕后掌握朝政枢纽”的前期动员与资源准备}

\eventanchor{WH-LEDGER-026-2}{前195年四月·“高祖刘邦去世，惠帝即位，吕后掌握朝政枢纽”的命令、联盟或地方响应}

\eventanchor{WH-LEDGER-026-3}{前195年四月·高祖刘邦去世，惠帝即位，吕后掌握朝政枢纽}

\eventanchor{WH-LEDGER-026-4}{前195年四月·“高祖刘邦去世，惠帝即位，吕后掌握朝政枢纽”的执行中的空间与人员处置}

\eventanchor{WH-LEDGER-026-5}{前195年四月·“高祖刘邦去世，惠帝即位，吕后掌握朝政枢纽”的后续制度或军政调整}

\eventanchor{WH-LEDGER-026-6}{前195年四月·“高祖刘邦去世，惠帝即位，吕后掌握朝政枢纽”的异源材料与影响范围的复核}

\eventanchor{WH-LEDGER-027-1}{前194年至前188年·“吕后在惠帝朝主导高层人事并处置戚夫人、赵王如意”的前期动员与资源准备}

\eventanchor{WH-LEDGER-027-2}{前194年至前188年·“吕后在惠帝朝主导高层人事并处置戚夫人、赵王如意”的命令、联盟或地方响应}

\eventanchor{WH-LEDGER-027-3}{前194年至前188年·吕后在惠帝朝主导高层人事并处置戚夫人、赵王如意}

\eventanchor{WH-LEDGER-027-4}{前194年至前188年·“吕后在惠帝朝主导高层人事并处置戚夫人、赵王如意”的执行中的空间与人员处置}

\eventanchor{WH-LEDGER-027-5}{前194年至前188年·“吕后在惠帝朝主导高层人事并处置戚夫人、赵王如意”的后续制度或军政调整}

\eventanchor{WH-LEDGER-027-6}{前194年至前188年·“吕后在惠帝朝主导高层人事并处置戚夫人、赵王如意”的异源材料与影响范围的复核}

\eventanchor{WH-LEDGER-028-1}{前188年八月·“惠帝去世，吕后临朝并先后立少帝”的前期动员与资源准备}

\eventanchor{WH-LEDGER-028-2}{前188年八月·“惠帝去世，吕后临朝并先后立少帝”的命令、联盟或地方响应}

\eventanchor{WH-LEDGER-028-3}{前188年八月·惠帝去世，吕后临朝并先后立少帝}

\eventanchor{WH-LEDGER-028-4}{前188年八月·“惠帝去世，吕后临朝并先后立少帝”的执行中的空间与人员处置}

\eventanchor{WH-LEDGER-028-5}{前188年八月·“惠帝去世，吕后临朝并先后立少帝”的后续制度或军政调整}

\eventanchor{WH-LEDGER-028-6}{前188年八月·“惠帝去世，吕后临朝并先后立少帝”的异源材料与影响范围的复核}

\eventanchor{WH-LEDGER-029-1}{前184年前后·“吕后封吕氏诸王并任用吕产、吕禄等掌军政”的前期动员与资源准备}

\eventanchor{WH-LEDGER-029-2}{前184年前后·“吕后封吕氏诸王并任用吕产、吕禄等掌军政”的命令、联盟或地方响应}

\eventanchor{WH-LEDGER-029-3}{前184年前后·吕后封吕氏诸王并任用吕产、吕禄等掌军政}

\eventanchor{WH-LEDGER-029-4}{前184年前后·“吕后封吕氏诸王并任用吕产、吕禄等掌军政”的执行中的空间与人员处置}

\eventanchor{WH-LEDGER-029-5}{前184年前后·“吕后封吕氏诸王并任用吕产、吕禄等掌军政”的后续制度或军政调整}

\eventanchor{WH-LEDGER-029-6}{前184年前后·“吕后封吕氏诸王并任用吕产、吕禄等掌军政”的异源材料与影响范围的复核}

\eventanchor{WH-LEDGER-030-1}{前180年八月·“吕后去世后，周勃、陈平等发动政变诛灭吕氏”的前期动员与资源准备}

\eventanchor{WH-LEDGER-030-2}{前180年八月·“吕后去世后，周勃、陈平等发动政变诛灭吕氏”的命令、联盟或地方响应}

\eventanchor{WH-LEDGER-030-3}{前180年八月·吕后去世后，周勃、陈平等发动政变诛灭吕氏}

\eventanchor{WH-LEDGER-030-4}{前180年八月·“吕后去世后，周勃、陈平等发动政变诛灭吕氏”的执行中的空间与人员处置}

\eventanchor{WH-LEDGER-030-5}{前180年八月·“吕后去世后，周勃、陈平等发动政变诛灭吕氏”的后续制度或军政调整}

\eventanchor{WH-LEDGER-030-6}{前180年八月·“吕后去世后，周勃、陈平等发动政变诛灭吕氏”的异源材料与影响范围的复核}

\eventanchor{WH-LEDGER-031-1}{前180年九月·“代王刘恒入长安即帝位，是为文帝”的前期动员与资源准备}

\eventanchor{WH-LEDGER-031-2}{前180年九月·“代王刘恒入长安即帝位，是为文帝”的命令、联盟或地方响应}

\eventanchor{WH-LEDGER-031-3}{前180年九月·代王刘恒入长安即帝位，是为文帝}

\eventanchor{WH-LEDGER-031-4}{前180年九月·“代王刘恒入长安即帝位，是为文帝”的执行中的空间与人员处置}

\eventanchor{WH-LEDGER-031-5}{前180年九月·“代王刘恒入长安即帝位，是为文帝”的后续制度或军政调整}

\eventanchor{WH-LEDGER-031-6}{前180年九月·“代王刘恒入长安即帝位，是为文帝”的异源材料与影响范围的复核}

\eventanchor{WH-LEDGER-032-1}{前178年·“文帝诏减田租为三十税一”的前期动员与资源准备}

\eventanchor{WH-LEDGER-032-2}{前178年·“文帝诏减田租为三十税一”的命令、联盟或地方响应}

\eventanchor{WH-LEDGER-032-3}{前178年·文帝诏减田租为三十税一}

\eventanchor{WH-LEDGER-032-4}{前178年·“文帝诏减田租为三十税一”的执行中的空间与人员处置}

\eventanchor{WH-LEDGER-032-5}{前178年·“文帝诏减田租为三十税一”的后续制度或军政调整}

\eventanchor{WH-LEDGER-032-6}{前178年·“文帝诏减田租为三十税一”的异源材料与影响范围的复核}

\eventanchor{WH-LEDGER-033-1}{前177年前后·“文帝弛山泽禁并许可民间铸钱，后引发货币质量问题”的前期动员与资源准备}

\eventanchor{WH-LEDGER-033-2}{前177年前后·“文帝弛山泽禁并许可民间铸钱，后引发货币质量问题”的命令、联盟或地方响应}

\eventanchor{WH-LEDGER-033-3}{前177年前后·文帝弛山泽禁并许可民间铸钱，后引发货币质量问题}

\eventanchor{WH-LEDGER-033-4}{前177年前后·“文帝弛山泽禁并许可民间铸钱，后引发货币质量问题”的执行中的空间与人员处置}

\eventanchor{WH-LEDGER-033-5}{前177年前后·“文帝弛山泽禁并许可民间铸钱，后引发货币质量问题”的后续制度或军政调整}

\eventanchor{WH-LEDGER-033-6}{前177年前后·“文帝弛山泽禁并许可民间铸钱，后引发货币质量问题”的异源材料与影响范围的复核}

\eventanchor{NS-LEDGER-008-3}{483年·齐武帝即位}

\subsection{事件导航与来源边界}

\eventanchor{NS-LEDGER-008-4}{483年·“齐武帝即位”的执行中的空间与人员处置}

\eventanchor{NS-LEDGER-008-5}{483年·“齐武帝即位”的后续制度或军政调整}

\eventanchor{NS-LEDGER-008-6}{483年·“齐武帝即位”的异源材料与影响范围的复核}

\eventanchor{NS-LEDGER-009-1}{494年·“齐明帝即位”的前期动员与资源准备}

\eventanchor{NS-LEDGER-009-2}{494年·“齐明帝即位”的命令、联盟或地方响应}

\eventanchor{NS-LEDGER-009-3}{494年·齐明帝即位}

\eventanchor{NS-LEDGER-009-4}{494年·“齐明帝即位”的执行中的空间与人员处置}

\eventanchor{NS-LEDGER-009-5}{494年·“齐明帝即位”的后续制度或军政调整}

\eventanchor{NS-LEDGER-009-6}{494年·“齐明帝即位”的异源材料与影响范围的复核}

\eventanchor{NS-LEDGER-010-1}{501年·“萧衍起兵”的前期动员与资源准备}

\eventanchor{NS-LEDGER-010-2}{501年·“萧衍起兵”的命令、联盟或地方响应}

\eventanchor{NS-LEDGER-010-3}{501年·萧衍起兵}

\eventanchor{NS-LEDGER-010-4}{501年·“萧衍起兵”的执行中的空间与人员处置}

\eventanchor{NS-LEDGER-010-5}{501年·“萧衍起兵”的后续制度或军政调整}

\eventanchor{NS-LEDGER-010-6}{501年·“萧衍起兵”的异源材料与影响范围的复核}

\eventanchor{NS-LEDGER-011-1}{502年·“萧衍建梁”的前期动员与资源准备}

\eventanchor{NS-LEDGER-011-2}{502年·“萧衍建梁”的命令、联盟或地方响应}

\eventanchor{NS-LEDGER-011-3}{502年·萧衍建梁}

\eventanchor{NS-LEDGER-011-4}{502年·“萧衍建梁”的执行中的空间与人员处置}

\eventanchor{NS-LEDGER-011-5}{502年·“萧衍建梁”的后续制度或军政调整}

\eventanchor{NS-LEDGER-011-6}{502年·“萧衍建梁”的异源材料与影响范围的复核}

\eventanchor{NS-LEDGER-012-1}{504年·“梁武帝公开崇佛政策”的前期动员与资源准备}

\eventanchor{NS-LEDGER-012-2}{504年·“梁武帝公开崇佛政策”的命令、联盟或地方响应}

\eventanchor{NS-LEDGER-012-3}{504年·梁武帝公开崇佛政策}

\eventanchor{NS-LEDGER-012-4}{504年·“梁武帝公开崇佛政策”的执行中的空间与人员处置}

\eventanchor{NS-LEDGER-012-5}{504年·“梁武帝公开崇佛政策”的后续制度或军政调整}

\eventanchor{NS-LEDGER-012-6}{504年·“梁武帝公开崇佛政策”的异源材料与影响范围的复核}

\eventanchor{NS-LEDGER-013-1}{523年·“六镇起义爆发”的前期动员与资源准备}

\eventanchor{NS-LEDGER-013-2}{523年·“六镇起义爆发”的命令、联盟或地方响应}

\eventanchor{NS-LEDGER-013-3}{523年·六镇起义爆发}

\eventanchor{NS-LEDGER-013-4}{523年·“六镇起义爆发”的执行中的空间与人员处置}

\eventanchor{NS-LEDGER-013-5}{523年·“六镇起义爆发”的后续制度或军政调整}

\eventanchor{NS-LEDGER-013-6}{523年·“六镇起义爆发”的异源材料与影响范围的复核}

\eventanchor{NS-LEDGER-014-1}{528年·“孝明帝死，尔朱荣入洛”的前期动员与资源准备}

\eventanchor{NS-LEDGER-014-2}{528年·“孝明帝死，尔朱荣入洛”的命令、联盟或地方响应}

\eventanchor{NS-LEDGER-014-3}{528年·孝明帝死，尔朱荣入洛}

\eventanchor{NS-LEDGER-014-4}{528年·“孝明帝死，尔朱荣入洛”的执行中的空间与人员处置}

\eventanchor{NS-LEDGER-014-5}{528年·“孝明帝死，尔朱荣入洛”的后续制度或军政调整}

\eventanchor{NS-LEDGER-014-6}{528年·“孝明帝死，尔朱荣入洛”的异源材料与影响范围的复核}

\eventanchor{NS-LEDGER-015-1}{530年·“孝庄帝杀尔朱荣”的前期动员与资源准备}

\eventanchor{NS-LEDGER-015-2}{530年·“孝庄帝杀尔朱荣”的命令、联盟或地方响应}

\eventanchor{NS-LEDGER-015-3}{530年·孝庄帝杀尔朱荣}

\eventanchor{NS-LEDGER-015-4}{530年·“孝庄帝杀尔朱荣”的执行中的空间与人员处置}

\eventanchor{NS-LEDGER-015-5}{530年·“孝庄帝杀尔朱荣”的后续制度或军政调整}

\eventanchor{NS-LEDGER-015-6}{530年·“孝庄帝杀尔朱荣”的异源材料与影响范围的复核}

\eventanchor{NS-LEDGER-016-1}{532年·“高欢控制东部朝廷”的前期动员与资源准备}

\eventanchor{NS-LEDGER-016-2}{532年·“高欢控制东部朝廷”的命令、联盟或地方响应}

\eventanchor{NS-LEDGER-016-3}{532年·高欢控制东部朝廷}

\eventanchor{NS-LEDGER-016-4}{532年·“高欢控制东部朝廷”的执行中的空间与人员处置}

\eventanchor{NS-LEDGER-016-5}{532年·“高欢控制东部朝廷”的后续制度或军政调整}

\eventanchor{NS-LEDGER-016-6}{532年·“高欢控制东部朝廷”的异源材料与影响范围的复核}

\eventanchor{NS-LEDGER-017-1}{534年·“北魏分裂为东魏西魏”的前期动员与资源准备}

\eventanchor{NS-LEDGER-017-2}{534年·“北魏分裂为东魏西魏”的命令、联盟或地方响应}

\eventanchor{NS-LEDGER-017-3}{534年·北魏分裂为东魏西魏}

\eventanchor{NS-LEDGER-017-4}{534年·“北魏分裂为东魏西魏”的执行中的空间与人员处置}

\eventanchor{NS-LEDGER-017-5}{534年·“北魏分裂为东魏西魏”的后续制度或军政调整}

\eventanchor{NS-LEDGER-017-6}{534年·“北魏分裂为东魏西魏”的异源材料与影响范围的复核}

\eventanchor{NS-LEDGER-018-1}{535年·“高欢经营邺城军政”的前期动员与资源准备}

\eventanchor{NS-LEDGER-018-2}{535年·“高欢经营邺城军政”的命令、联盟或地方响应}

\eventanchor{NS-LEDGER-018-3}{535年·高欢经营邺城军政}

\subsection{事件导航与来源边界}

\eventanchor{NS-LEDGER-018-4}{535年·“高欢经营邺城军政”的执行中的空间与人员处置}

\eventanchor{NS-LEDGER-018-5}{535年·“高欢经营邺城军政”的后续制度或军政调整}

\eventanchor{NS-LEDGER-018-6}{535年·“高欢经营邺城军政”的异源材料与影响范围的复核}

\eventanchor{NS-LEDGER-019-1}{535年·“宇文泰控制长安”的前期动员与资源准备}

\eventanchor{NS-LEDGER-019-2}{535年·“宇文泰控制长安”的命令、联盟或地方响应}

\eventanchor{NS-LEDGER-019-3}{535年·宇文泰控制长安}

\eventanchor{NS-LEDGER-019-4}{535年·“宇文泰控制长安”的执行中的空间与人员处置}

\eventanchor{NS-LEDGER-019-5}{535年·“宇文泰控制长安”的后续制度或军政调整}

\eventanchor{NS-LEDGER-019-6}{535年·“宇文泰控制长安”的异源材料与影响范围的复核}

\eventanchor{NS-LEDGER-020-1}{547年·“侯景叛高欢南投梁”的前期动员与资源准备}

\eventanchor{NS-LEDGER-020-2}{547年·“侯景叛高欢南投梁”的命令、联盟或地方响应}

\eventanchor{NS-LEDGER-020-3}{547年·侯景叛高欢南投梁}

\eventanchor{NS-LEDGER-020-4}{547年·“侯景叛高欢南投梁”的执行中的空间与人员处置}

\eventanchor{NS-LEDGER-020-5}{547年·“侯景叛高欢南投梁”的后续制度或军政调整}

\eventanchor{NS-LEDGER-020-6}{547年·“侯景叛高欢南投梁”的异源材料与影响范围的复核}

\eventanchor{NS-LEDGER-021-1}{548年·“侯景渡江反梁”的前期动员与资源准备}

\eventanchor{NS-LEDGER-021-2}{548年·“侯景渡江反梁”的命令、联盟或地方响应}

\eventanchor{NS-LEDGER-021-3}{548年·侯景渡江反梁}

\eventanchor{NS-LEDGER-021-4}{548年·“侯景渡江反梁”的执行中的空间与人员处置}

\eventanchor{NS-LEDGER-021-5}{548年·“侯景渡江反梁”的后续制度或军政调整}

\eventanchor{NS-LEDGER-021-6}{548年·“侯景渡江反梁”的异源材料与影响范围的复核}

\eventanchor{NS-LEDGER-022-1}{549年·“侯景攻陷建康”的前期动员与资源准备}

\eventanchor{NS-LEDGER-022-2}{549年·“侯景攻陷建康”的命令、联盟或地方响应}

\eventanchor{NS-LEDGER-022-3}{549年·侯景攻陷建康}

\eventanchor{NS-LEDGER-022-4}{549年·“侯景攻陷建康”的执行中的空间与人员处置}

\eventanchor{NS-LEDGER-022-5}{549年·“侯景攻陷建康”的后续制度或军政调整}

\eventanchor{NS-LEDGER-022-6}{549年·“侯景攻陷建康”的异源材料与影响范围的复核}

\eventanchor{NS-LEDGER-023-1}{552年·“王僧辩陈霸先平侯景”的前期动员与资源准备}

\eventanchor{NS-LEDGER-023-2}{552年·“王僧辩陈霸先平侯景”的命令、联盟或地方响应}

\eventanchor{NS-LEDGER-023-3}{552年·王僧辩陈霸先平侯景}

\eventanchor{NS-LEDGER-023-4}{552年·“王僧辩陈霸先平侯景”的执行中的空间与人员处置}

\eventanchor{NS-LEDGER-023-5}{552年·“王僧辩陈霸先平侯景”的后续制度或军政调整}

\eventanchor{NS-LEDGER-023-6}{552年·“王僧辩陈霸先平侯景”的异源材料与影响范围的复核}

\eventanchor{NS-LEDGER-024-1}{554年·“西魏攻取江陵”的前期动员与资源准备}

\eventanchor{NS-LEDGER-024-2}{554年·“西魏攻取江陵”的命令、联盟或地方响应}

\eventanchor{NS-LEDGER-024-3}{554年·西魏攻取江陵}

\eventanchor{NS-LEDGER-024-4}{554年·“西魏攻取江陵”的执行中的空间与人员处置}

\eventanchor{NS-LEDGER-024-5}{554年·“西魏攻取江陵”的后续制度或军政调整}

\eventanchor{NS-LEDGER-024-6}{554年·“西魏攻取江陵”的异源材料与影响范围的复核}

\eventanchor{NS-LEDGER-025-1}{557年·“陈霸先建陈”的前期动员与资源准备}

\eventanchor{NS-LEDGER-025-2}{557年·“陈霸先建陈”的命令、联盟或地方响应}

\eventanchor{NS-LEDGER-025-3}{557年·陈霸先建陈}

\eventanchor{NS-LEDGER-025-4}{557年·“陈霸先建陈”的执行中的空间与人员处置}

\eventanchor{NS-LEDGER-025-5}{557年·“陈霸先建陈”的后续制度或军政调整}

\eventanchor{NS-LEDGER-025-6}{557年·“陈霸先建陈”的异源材料与影响范围的复核}

\eventanchor{NS-LEDGER-026-1}{557年·“宇文觉建北周”的前期动员与资源准备}

\eventanchor{NS-LEDGER-026-2}{557年·“宇文觉建北周”的命令、联盟或地方响应}

\eventanchor{NS-LEDGER-026-3}{557年·宇文觉建北周}

\eventanchor{NS-LEDGER-026-4}{557年·“宇文觉建北周”的执行中的空间与人员处置}

\eventanchor{NS-LEDGER-026-5}{557年·“宇文觉建北周”的后续制度或军政调整}

\eventanchor{NS-LEDGER-026-6}{557年·“宇文觉建北周”的异源材料与影响范围的复核}

\eventanchor{NS-LEDGER-027-1}{565年·“后主继位”的前期动员与资源准备}

\eventanchor{NS-LEDGER-027-2}{565年·“后主继位”的命令、联盟或地方响应}

\eventanchor{NS-LEDGER-027-3}{565年·后主继位}

\eventanchor{NS-LEDGER-027-4}{565年·“后主继位”的执行中的空间与人员处置}

\eventanchor{NS-LEDGER-027-5}{565年·“后主继位”的后续制度或军政调整}

\eventanchor{NS-LEDGER-027-6}{565年·“后主继位”的异源材料与影响范围的复核}

\eventanchor{NS-LEDGER-028-1}{572年·“周武帝诛宇文护”的前期动员与资源准备}

\eventanchor{NS-LEDGER-028-2}{572年·“周武帝诛宇文护”的命令、联盟或地方响应}

\eventanchor{NS-LEDGER-028-3}{572年·周武帝诛宇文护}

\subsection{事件导航与来源边界}

\eventanchor{NS-LEDGER-028-4}{572年·“周武帝诛宇文护”的执行中的空间与人员处置}

\eventanchor{NS-LEDGER-028-5}{572年·“周武帝诛宇文护”的后续制度或军政调整}

\eventanchor{NS-LEDGER-028-6}{572年·“周武帝诛宇文护”的异源材料与影响范围的复核}

\eventanchor{NS-LEDGER-029-1}{574年·“周武帝禁佛道”的前期动员与资源准备}

\eventanchor{NS-LEDGER-029-2}{574年·“周武帝禁佛道”的命令、联盟或地方响应}

\eventanchor{NS-LEDGER-029-3}{574年·周武帝禁佛道}

\eventanchor{NS-LEDGER-029-4}{574年·“周武帝禁佛道”的执行中的空间与人员处置}

\eventanchor{NS-LEDGER-029-5}{574年·“周武帝禁佛道”的后续制度或军政调整}

\eventanchor{NS-LEDGER-029-6}{574年·“周武帝禁佛道”的异源材料与影响范围的复核}

\eventanchor{NS-LEDGER-030-1}{577年·“北周灭北齐”的前期动员与资源准备}

\eventanchor{NS-LEDGER-030-2}{577年·“北周灭北齐”的命令、联盟或地方响应}

\eventanchor{NS-LEDGER-030-3}{577年·北周灭北齐}

\eventanchor{NS-LEDGER-030-4}{577年·“北周灭北齐”的执行中的空间与人员处置}

\eventanchor{NS-LEDGER-030-5}{577年·“北周灭北齐”的后续制度或军政调整}

\eventanchor{NS-LEDGER-030-6}{577年·“北周灭北齐”的异源材料与影响范围的复核}

\eventanchor{NS-LEDGER-031-1}{578年·“周武帝去世”的前期动员与资源准备}

\eventanchor{NS-LEDGER-031-2}{578年·“周武帝去世”的命令、联盟或地方响应}

\eventanchor{NS-LEDGER-031-3}{578年·周武帝去世}

\eventanchor{NS-LEDGER-031-4}{578年·“周武帝去世”的执行中的空间与人员处置}

\eventanchor{NS-LEDGER-031-5}{578年·“周武帝去世”的后续制度或军政调整}

\eventanchor{NS-LEDGER-031-6}{578年·“周武帝去世”的异源材料与影响范围的复核}

\eventanchor{NS-LEDGER-032-1}{580年·“杨坚辅政”的前期动员与资源准备}

\eventanchor{NS-LEDGER-032-2}{580年·“杨坚辅政”的命令、联盟或地方响应}

\eventanchor{NS-LEDGER-032-3}{580年·杨坚辅政}

\eventanchor{NS-LEDGER-032-4}{580年·“杨坚辅政”的执行中的空间与人员处置}

\eventanchor{NS-LEDGER-032-5}{580年·“杨坚辅政”的后续制度或军政调整}

\eventanchor{NS-LEDGER-032-6}{580年·“杨坚辅政”的异源材料与影响范围的复核}

\eventanchor{NS-LEDGER-033-1}{581年·“杨坚受禅建隋”的前期动员与资源准备}

\eventanchor{NS-LEDGER-033-2}{581年·“杨坚受禅建隋”的命令、联盟或地方响应}

\eventanchor{NS-LEDGER-033-3}{581年·杨坚受禅建隋}

\eventanchor{NS-LEDGER-033-4}{581年·“杨坚受禅建隋”的执行中的空间与人员处置}

\eventanchor{NS-LEDGER-033-5}{581年·“杨坚受禅建隋”的后续制度或军政调整}

\eventanchor{NS-LEDGER-033-6}{581年·“杨坚受禅建隋”的异源材料与影响范围的复核}

\eventanchor{NS-LEDGER-034-1}{583年·“隋整顿北方军政与户籍”的前期动员与资源准备}

\eventanchor{NS-LEDGER-034-2}{583年·“隋整顿北方军政与户籍”的命令、联盟或地方响应}

\eventanchor{NS-LEDGER-034-3}{583年·隋整顿北方军政与户籍}

\eventanchor{NS-LEDGER-034-4}{583年·“隋整顿北方军政与户籍”的执行中的空间与人员处置}

\eventanchor{NS-LEDGER-034-5}{583年·“隋整顿北方军政与户籍”的后续制度或军政调整}

\eventanchor{NS-LEDGER-034-6}{583年·“隋整顿北方军政与户籍”的异源材料与影响范围的复核}

\eventanchor{NS-LEDGER-035-1}{588年·“隋军南下伐陈”的前期动员与资源准备}

\eventanchor{NS-LEDGER-035-2}{588年·“隋军南下伐陈”的命令、联盟或地方响应}

\eventanchor{NS-LEDGER-035-3}{588年·隋军南下伐陈}

\eventanchor{NS-LEDGER-035-4}{588年·“隋军南下伐陈”的执行中的空间与人员处置}

\eventanchor{NS-LEDGER-035-5}{588年·“隋军南下伐陈”的后续制度或军政调整}

\eventanchor{NS-LEDGER-035-6}{588年·“隋军南下伐陈”的异源材料与影响范围的复核}

\eventanchor{NS-LEDGER-036-1}{589年·“建康失陷，陈亡”的前期动员与资源准备}

\eventanchor{NS-LEDGER-036-2}{589年·“建康失陷，陈亡”的命令、联盟或地方响应}

\eventanchor{NS-LEDGER-036-3}{589年·建康失陷，陈亡}

\eventanchor{NS-LEDGER-036-4}{589年·“建康失陷，陈亡”的执行中的空间与人员处置}

\eventanchor{NS-LEDGER-036-5}{589年·“建康失陷，陈亡”的后续制度或军政调整}

\eventanchor{NS-LEDGER-036-6}{589年·“建康失陷，陈亡”的异源材料与影响范围的复核}

\eventanchor{NS-LEDGER-037-1}{450年·“北魏南侵至江淮”的前期动员与资源准备}

\eventanchor{NS-LEDGER-037-2}{450年·“北魏南侵至江淮”的命令、联盟或地方响应}

\eventanchor{NS-LEDGER-037-3}{450年·北魏南侵至江淮}

\eventanchor{NS-LEDGER-037-4}{450年·“北魏南侵至江淮”的执行中的空间与人员处置}

\eventanchor{NS-LEDGER-037-5}{450年·“北魏南侵至江淮”的后续制度或军政调整}

\eventanchor{NS-LEDGER-037-6}{450年·“北魏南侵至江淮”的异源材料与影响范围的复核}

\eventanchor{NS-LEDGER-038-1}{471年·“献文帝禅位，孝文帝继位”的前期动员与资源准备}

\eventanchor{NS-LEDGER-038-2}{471年·“献文帝禅位，孝文帝继位”的命令、联盟或地方响应}

\eventanchor{NS-LEDGER-038-3}{471年·献文帝禅位，孝文帝继位}

\subsection{事件导航与来源边界}

\eventanchor{NS-LEDGER-038-4}{471年·“献文帝禅位，孝文帝继位”的执行中的空间与人员处置}

\eventanchor{NS-LEDGER-038-5}{471年·“献文帝禅位，孝文帝继位”的后续制度或军政调整}

\eventanchor{NS-LEDGER-038-6}{471年·“献文帝禅位，孝文帝继位”的异源材料与影响范围的复核}

\eventanchor{NS-LEDGER-039-1}{494年·“北魏迁都洛阳”的前期动员与资源准备}

\eventanchor{NS-LEDGER-039-2}{494年·“北魏迁都洛阳”的命令、联盟或地方响应}

\eventanchor{NS-LEDGER-039-3}{494年·北魏迁都洛阳}

\eventanchor{NS-LEDGER-039-4}{494年·“北魏迁都洛阳”的执行中的空间与人员处置}

\eventanchor{NS-LEDGER-039-5}{494年·“北魏迁都洛阳”的后续制度或军政调整}

\eventanchor{NS-LEDGER-039-6}{494年·“北魏迁都洛阳”的异源材料与影响范围的复核}

\eventanchor{NS-LEDGER-040-1}{495年·“孝文帝改革姓氏服制”的前期动员与资源准备}

\eventanchor{NS-LEDGER-040-2}{495年·“孝文帝改革姓氏服制”的命令、联盟或地方响应}

\eventanchor{NS-LEDGER-040-3}{495年·孝文帝改革姓氏服制}

\eventanchor{NS-LEDGER-040-4}{495年·“孝文帝改革姓氏服制”的执行中的空间与人员处置}

\eventanchor{NS-LEDGER-040-5}{495年·“孝文帝改革姓氏服制”的后续制度或军政调整}

\eventanchor{NS-LEDGER-040-6}{495年·“孝文帝改革姓氏服制”的异源材料与影响范围的复核}

\eventanchor{NSL-440-589-265}{579年·周宣帝禅位于幼子静帝，自为天元皇帝}

\eventanchor{NSL-440-589-266}{579年·北周在河北推行州郡整并、户籍和赋役恢复措施}

\eventanchor{NSL-440-589-267}{579年·北周对佛教禁令有所松动，部分寺院和僧侣重新活动}

\eventanchor{NSL-440-589-268}{580年·周宣帝去世，杨坚以外戚身份辅政静帝}

\eventanchor{NSL-440-589-269}{580年·尉迟迥在相州起兵反杨坚，河北旧齐势力与部分军将响应}

\eventanchor{NSL-440-589-270}{580年·杨坚平定尉迟迥后重新任命河北州郡和军队将领}

\eventanchor{NSL-440-589-271}{581年·杨坚受禅建立隋，北周灭亡}

\eventanchor{NSL-440-589-272}{581年·隋初重整中央官制、府兵与州郡任命，削减北周宗室影响}

\eventanchor{NSL-440-589-273}{581年·突厥沙钵略可汗支持北周宗室反隋，北部边境紧张}

\eventanchor{NSL-440-589-274}{582年·陈宣帝去世，陈叔宝即位，是为后主}

\eventanchor{NSL-440-589-275}{582年·隋营建大兴城，规划宫城、官署、坊市与交通供给}

\eventanchor{NSL-440-589-276}{582年·突厥大举南下，隋以长城、州镇和机动军队应对}

\eventanchor{NSL-440-589-277}{583年·隋修订开皇律令，强调简约刑法与中央审断}

\eventanchor{NSL-440-589-278}{583年·隋整并州郡、清理户籍与赋役，推进均田和租调体系}

\eventanchor{NSL-440-589-279}{583年·隋利用突厥可汗间竞争，削弱沙钵略对北周遗臣的支持}

\eventanchor{NSL-440-589-280}{584年·隋开通广通渠，连接大兴城与潼关、黄河运输体系}

\eventanchor{NSL-440-589-281}{584年·陈朝面对隋北方整合，修整长江防线并调整江北守将}

\eventanchor{NSL-440-589-282}{585年·隋在北方旧齐地区继续编户、核田与恢复仓储转运}

\eventanchor{NSL-440-589-283}{585年·陈后主朝廷的奢侈、宫廷文学与近侍政治受到史家强烈批评}

\eventanchor{NSL-440-589-284}{586年·隋继续整顿府兵、仓储和江淮方向的军情侦察}

\eventanchor{NSL-440-589-285}{586年·陈朝在长江下游维持水军、城戍与漕运，承受持续军费}

\eventanchor{NSL-440-589-286}{587年·隋废后梁，接收荆襄州郡与江陵资源}

\eventanchor{NSL-440-589-287}{587年·隋在荆襄安置后梁官员、军人和居民，整顿道路与粮运}

\eventanchor{NSL-440-589-288}{588年·隋以杨广、杨俊、杨素等分路部署，发动灭陈战争}

\eventanchor{NSL-440-589-289}{588年·陈朝调集长江水军和沿岸守将抵御隋军}

\eventanchor{NSL-440-589-290}{589年·隋军渡江攻入建康，陈后主被俘，陈亡}

\eventanchor{NSL-440-589-291}{589年·隋接收陈朝州郡、户籍、仓储和官员，派使安抚江南}

\eventanchor{NSL-440-589-292}{589年·隋处理陈室宗族、降将与江南士族的封授和迁置}

\eventanchor{NSL-440-589-293}{约589年·颜之推晚年完成或修订《颜氏家训》的主要篇章，反思北齐、北周、隋之间的家族处境}

\eventanchor{NSL-440-589-294}{约550年·北齐在邺城重整东魏旧都的宫署、市场与军政居住空间}

\eventanchor{NSL-440-589-295}{552年·北齐在北部边地处置柔然覆亡后逃入境内的部众}

\eventanchor{NSL-440-589-296}{554年·江陵陷落后，大批梁朝官员、工匠和图书被迁徙或散失}

\eventanchor{NSL-440-589-297}{557年·北周建国初年以六官名目协调军政、财政与礼制机构}

\eventanchor{NSL-440-589-298}{559年·陈朝在建康战后修复宫城、仓储和长江水运节点}

\eventanchor{NSL-440-589-299}{563年·北齐河清制度整理涉及律令、官品与地方行政的重新表述}

\eventanchor{NSL-440-589-300}{566年·北齐继续以贡赐和使节维系对突厥的边境关系}

\eventanchor{NSL-440-589-301}{570年·北周侦察并经营河东、河南道路，准备对北齐发动更大攻势}

\eventanchor{NSL-440-589-302}{575年·禁佛后的北周以州郡接收寺院人户与土地，重估其赋役归属}

\eventanchor{NSL-440-589-303}{577年·北周在新并北齐地区重建仓储、道路与军队粮运}

\eventanchor{NSL-440-589-304}{580年·尉迟迥之乱中，杨坚军利用关中与河北粮道协调平叛}

\eventanchor{NSL-440-589-305}{582年·大兴城建设同步规划水源、道路与坊市，吸纳关中工匠和劳力}
